\documentclass[twocolumn]{aastex631}

\usepackage{amsmath}
\usepackage{hyperref}
\usepackage{xcolor}

\shorttitle{EPS Research Astro-RAG Platform}
\shortauthors{Flynn}

\begin{document}

\title{The EPS Research Astro-RAG Platform: A Unified Open-Science Infrastructure\\
for Cross-Epoch Astrophysical Kinematic Analysis, LLM-Assisted Research Workflows,\\
and Educational Outreach}

\author{David C. Flynn}
\affiliation{EPS Research, Laurel, MD 20707, USA}
\email{davidflynn@eps-research.com}
\href{https://orcid.org/0000-0002-2768-6650}{0000-0002-2768-6650}

\begin{abstract}
We present the EPS Research Astro-RAG Platform (v1.0), an open-science research
infrastructure providing four machine-readable, cross-epoch astrophysical corpora,
120 verified executable Jupyter notebook examples, a QuickStart reproducibility
pathway, an educational High-School Exploration Track, and a roadmap for
LLM-assisted retrieval-augmented generation (RAG) research workflows. The platform
spans the full observable cosmic epoch from the local universe ($z = 0$) to the
epoch approaching reionization ($z \sim 6$), providing the first unified
kinematic dataset connecting HI 21cm rotation curves, Milky Way globular cluster
dynamics, and high-redshift ALMA [CII] morpho-kinematics under a common
schema and analysis framework. The four corpora contain 772 total objects:
438 HI-traced galaxies (Unified HI Rotation Curve Corpus v7.0), 129
dwarf/irregular galaxies (Dwarf/Irregular HI Corpus v1.0), 174 Milky Way
globular clusters (GC Corpus v1.3.1), and 31 star-forming galaxies at
$z = 4.26$--$5.68$ (High-z Kinematic Corpus Z1). The Unified HI corpus
preserves the full fidelity of the Case Western Reserve University SPARC
database (Lelli et al. 2016), including all photometric decomposition components
($V_{\rm gas}$, $V_{\rm disk}$, $V_{\rm bul}$) and surface brightness profiles,
augmented with THINGS, LITTLE THINGS, and WALLABY DR2 data. The unifying
scientific framework is the omega kinematic correction of Flynn \& Cannaliato
(2025), which reveals a robust sign reversal in the boundary-point angular
velocity $\omega$ from positive values at $z = 0$ (SPARC mean $+7.06 \pm 3.26$
rad~Gyr$^{-1}$) to negative values at $z \sim 5$ (Z1 median $-13.05$
rad~Gyr$^{-1}$), motivating future RAMSES cosmological simulations. All
corpora are released under CC BY 4.0 at Zenodo with permanent DOIs, and the
platform is archived at \url{https://github.com/eps-research/rag-corpus-series}
(DOI: \href{https://doi.org/10.5281/zenodo.20398430}{10.5281/zenodo.20398430}).
\end{abstract}

\keywords{catalogs --- methods: data analysis --- galaxies: kinematics and dynamics --- galaxies: dwarf --- globular clusters: general --- galaxies: high-redshift}

%=====================================================================
\section{Introduction}
\label{sec:intro}
%=====================================================================

The intersection of large-scale astrophysical surveys and machine-learning
infrastructure has created both an opportunity and a challenge for the research
community. Surveys such as SPARC \citep{lelli2016}, THINGS \citep{walter2008},
LITTLE THINGS \citep{oh2015}, WALLABY \citep{koribalski2020}, and ALPINE
\citep{lefevre2020} have produced rich kinematic datasets spanning orders of
magnitude in galaxy mass, morphology, and cosmic epoch. However, these datasets
exist in heterogeneous formats across dozens of source papers, making systematic
cross-survey and cross-epoch analysis difficult for individual researchers.

The emergence of large language models (LLMs) and retrieval-augmented generation
(RAG) pipelines as scientific tools adds a further requirement: datasets must be
not only machine-readable but \emph{self-describing} --- structured such that an
LLM can ingest, query, and reason over the data without external preprocessing.
Standard FITS tables and survey-specific ASCII formats do not meet this
requirement.

We present the EPS Research Astro-RAG Platform, a unified open-science
infrastructure that addresses both challenges simultaneously. The platform
provides four corpora in a common schema (JSON + flat CSV + RAG-ready JSONL +
per-object ZIP), 120 verified executable Jupyter notebooks spanning all four
datasets, a reproducible QuickStart pathway, and a High-School Exploration
Track designed to make astrophysical research accessible to students without
prior domain knowledge.

The unifying scientific contribution across all four corpora is the omega
kinematic correction \citep{flynn2025}, which introduces a boundary-point
angular velocity parameter $\omega$ derived from the innermost and outermost
resolved rings of each rotation curve. Applied across 84 SPARC quality-1
galaxies, the correction yields a mean $\omega = +7.06 \pm 3.26$ rad~Gyr$^{-1}$
at $z = 0$. Extended to 8 confirmed rotators at $z \sim 4$--$6$ via ALMA
[CII] 3DBarolo models \citep{jones2021}, the same parameter yields a median
$\omega = -13.05$ rad~Gyr$^{-1}$ --- a sign reversal of $\sim 20$ rad~Gyr$^{-1}$
across $\sim 9$ Gyr of cosmic evolution.

This paper describes the platform architecture (Section~\ref{sec:architecture}),
the four corpora (Section~\ref{sec:corpora}), the example library and educational
track (Section~\ref{sec:examples}), the reproducibility infrastructure
(Section~\ref{sec:reproducibility}), and the planned LLM fine-tune team
(Section~\ref{sec:llms}). Data access and citation information are provided
in Section~\ref{sec:data}.

%=====================================================================
\section{Platform Architecture}
\label{sec:architecture}
%=====================================================================

The platform is organized into five functional silos, reflecting the natural
progression from raw data to scientific analysis to publication to AI-assisted
workflows:

\begin{enumerate}
\item \textbf{Silo 1 --- $z = 0$ Data}: Three local-universe corpora (HI galaxies,
      dwarf/irregular galaxies, Milky Way globular clusters).
\item \textbf{Silo 2 --- $z \sim 6$ Data}: The High-z Kinematic Corpus Z1,
      providing ALPINE [CII] morpho-kinematics at $z = 4.26$--$5.68$.
\item \textbf{Silo 3 --- Example Library}: 120 verified Jupyter notebooks
      plus the High-School Exploration Track.
\item \textbf{Silo 4 --- Papers \& Preprints}: The EPS Research publication
      arc with DOIs and arXiv identifiers.
\item \textbf{Silo 5 --- LLMs \& Tools}: Planned fine-tuned model adapters
      and RAG utilities (Stage 2).
\end{enumerate}

\subsection{Common Schema Design}
\label{sec:schema}

All four corpora share a common design philosophy. Each corpus is distributed in
four formats:

\begin{itemize}
\item \textbf{JSON}: Hierarchical, per-object records with full metadata and
      per-measurement data arrays. Self-describing field names with units embedded.
\item \textbf{Flat CSV}: One row per object, suitable for pandas, R, and
      spreadsheet analysis.
\item \textbf{JSONL}: One JSON object per line, optimized for LLM RAG pipelines
      and streaming ingestion.
\item \textbf{Per-object ZIP}: Individual JSON files for each object, enabling
      targeted retrieval without loading the full corpus.
\end{itemize}

Quality tiers are explicit in all corpora. For the HI and dwarf corpora,
Tier~1 objects have full per-ring kinematic data; Tier~2 objects have
integrated parameters only. For the Z1 corpus, Tier~1 objects are confirmed
3DBarolo rotators with per-ring $V_{\rm rot}$ and $\sigma$ profiles;
Tier~2 objects have morpho-kinematic classifications only.

%=====================================================================
\section{The Four Corpora}
\label{sec:corpora}
%=====================================================================

\subsection{Unified HI Rotation Curve Corpus v7.0}
\label{sec:hi}

The Unified HI Rotation Curve Corpus v7.0 \citep{flynn2026_hi} contains 438
spatially resolved rotation curves from four major HI surveys: SPARC
\citep[175 galaxies;][]{lelli2016}, THINGS \citep[34 galaxies;][]{walter2008},
LITTLE THINGS \citep[26 galaxies;][]{oh2015}, and WALLABY DR2
\citep[203 galaxies;][]{koribalski2020}.

\textbf{SPARC full-fidelity preservation.} The SPARC component of this corpus
preserves the complete photometric decomposition provided by the Case Western
Reserve University SPARC database \citep{lelli2016}. For each of the 175 SPARC
galaxies, the corpus retains all per-ring components: the observed rotation
velocity $V_{\rm obs}$ with error $\delta V$, the gas contribution $V_{\rm gas}$,
the stellar disk contribution $V_{\rm disk}$, the bulge contribution $V_{\rm bul}$
(where applicable), and the 3.6~$\mu$m (Spitzer IRAC) surface brightness profiles $\Sigma_{\rm disk}$
and $\Sigma_{\rm bul}$. This full-fidelity preservation enables baryonic mass
decomposition analyses and Tully--Fisher relation studies without requiring
access to the original survey database. The SPARC data have been cross-validated
against the published SPARC tables (Lelli et al. 2016, Table~1 and online
appendix) at the individual ring level.

The corpus includes 84 SPARC galaxies with the original Lelli et al.\ (2016) Q=1 quality flag (the ``golden sample'' used in Flynn \& Cannaliato 2025), identified by the \texttt{sparc\_quality=1} flag in the corpus schema.

\subsection{Dwarf/Irregular HI Corpus v1.0}
\label{sec:dwarfs}

The Dwarf/Irregular HI Corpus v1.0 \citep{flynn2026_dwarfs} contains 129
dwarf and irregular galaxies from LVHIS \citep{koribalski2018}, VLA-ANGST
\citep{ott2012}, LITTLE THINGS \citep{oh2015}, and WALLABY DR2
\citep{koribalski2020}. The corpus focuses on the low-mass, dark-matter-dominated
end of the galaxy mass function, providing a complement to the SPARC-dominated
HI corpus.

Of 129 total galaxies, 24 pass the ``omega-ready'' criterion (quality tier~1,
at least 5 resolved rings, and kinematically regular morphology), yielding a dwarf
corpus median $\omega = +9.94$ rad~Gyr$^{-1}$ --- higher than the SPARC mean,
consistent with the expectation that dark-matter-dominated systems with
rising outer rotation curves produce larger positive omega values.

\subsection{Milky Way Globular Cluster Corpus v1.3.1}
\label{sec:gc}

The Milky Way Globular Cluster Corpus v1.3.1 \citep{flynn2026_gc} contains
174 Milky Way globular clusters assembled from four primary sources: the
Harris (1996, 2010 edition) photometric catalog \citep{harris1996}, Vasiliev
\& Baumgardt (2021) Gaia EDR3 proper motions \citep{vasiliev2021}, HST structural
parameters, and Baumgardt \& Hilker (2018) N-body dynamical models
\citep{baumgardt2018}. All four source catalogs are cross-matched by cluster
identifier and harmonized into a single per-cluster JSON record.

\subsection{High-z Kinematic Corpus Z1}
\label{sec:z1}

The High-z Kinematic Corpus Z1 \citep{flynn2026_z1} contains 31 ALPINE
\citep{lefevre2020} star-forming galaxies at $z = 4.26$--$5.68$, drawn from
the Jones et al. (2021) morpho-kinematic classification catalog
\citep{jones2021}. Eight galaxies are confirmed Tier~1 rotators with per-ring
3DBarolo rotation curves; 23 additional galaxies have morpho-kinematic
classifications only (Tier~2).

The corpus provides the high-redshift anchor of the platform, enabling the
first direct cross-epoch omega comparison between local HI rotation curves
and ALMA [CII] kinematics near the epoch of reionization.

\textbf{Cross-epoch result.} All 8 Tier~1 Z1 rotators show negative omega
values (median $-13.05$ rad~Gyr$^{-1}$) under the Flynn \& Cannaliato (2025)
kinematic correction, contrasting with positive values at $z = 0$ (SPARC mean
$+7.06 \pm 3.26$ rad~Gyr$^{-1}$; dwarf median $+9.94$ rad~Gyr$^{-1}$). This
sign reversal of $\sim 20$ rad~Gyr$^{-1}$ across $\sim 9$ Gyr of cosmic
evolution is consistent with the structural transformation from compact,
centrally concentrated high-$z$ star-forming systems to extended rotating
disks at $z = 0$. RAMSES cosmological simulations tracing this evolution from
$z = 6$ initial conditions are planned for 2026--2027.

%=====================================================================
\section{Example Library and Educational Track}
\label{sec:examples}
%=====================================================================

\subsection{120 Verified Jupyter Notebooks}
\label{sec:notebooks}

The platform provides 120 executable Jupyter notebooks organized into five
groups, all verified to execute without errors via automated \texttt{nbconvert}
end-to-end testing:

\begin{itemize}
\item \textbf{SPARC/HI Examples} (25 notebooks): Rotation curve plotting,
      baryonic decomposition, omega correction, WALLABY Tier~2 analysis.
\item \textbf{Dwarf/Irregular Examples} (25 notebooks): Omega-ready galaxy
      selection, DDO154/DDO161 cross-analysis, LVHIS and VLA-ANGST comparisons.
\item \textbf{Globular Cluster Examples} (25 notebooks): Proper motion queries,
      N-body mass modeling, APOGEE chemistry cross-matching, multi-survey
      parameter comparison.
\item \textbf{High-z Examples} (25 notebooks): [CII] rotation curve analysis,
      ALPINE population statistics, W15 disk criteria evaluation, cross-corpus
      omega bridge.
\item \textbf{High-School Exploration Track} (20 notebooks): Two-track
      curriculum (Track A: ages 12--14; Track B: ages 15--18) introducing
      galaxy kinematics, dark matter, and the omega correction without
      requiring prior astrophysics background.
\end{itemize}

All notebooks require only \texttt{numpy}, \texttt{matplotlib}, and
\texttt{jupyter} as dependencies, with \texttt{scipy} for statistical notebooks.
No proprietary software, telescope archive accounts, or institutional access
is required.

\subsection{High-School Exploration Track}
\label{sec:hs}

The High-School Exploration Track is a novel feature that distinguishes the
EPS Research platform from standard data releases. Track A (10 notebooks) is
designed for students aged 12--14 with no prior astrophysics knowledge,
introducing concepts such as galaxy morphology, rotation curves, and dark matter
through guided data exploration. Track B (10 notebooks) targets ages 15--18
with a calculus-ready audience, introducing the omega correction formula,
RMSE as a model evaluation metric, and cross-epoch kinematics.

Both tracks load directly from the same Zenodo corpus files used by the
research notebooks, ensuring that students interact with real, published
scientific data rather than simplified approximations. Cloud access via
Google Colab (planned Stage 2) will remove the hardware barrier entirely,
enabling in-classroom use without requiring students to install software.

%=====================================================================
\section{Reproducibility Infrastructure}
\label{sec:reproducibility}
%=====================================================================

\subsection{QuickStart Pathway}
\label{sec:quickstart}

The platform provides a \texttt{QuickStart.ipynb} notebook at the repository
root that implements a ``golden path'' through all four corpora. Starting from
a fresh clone, a user can reproduce the core cross-epoch omega sign reversal
result in under 10 minutes using only:

\begin{verbatim}
git clone https://github.com/eps-research/
         rag-corpus-series
cd rag-corpus-series
pip install -r requirements.txt
python download_corpora.py
jupyter lab
\end{verbatim}

The \texttt{download\_corpora.py} script fetches all four corpus JSON files
directly from their Zenodo DOIs and places them in the correct directories
for both the QuickStart notebook and all 120 example notebooks.

\subsection{Verification Protocol}
\label{sec:verification}

All 120 notebooks were verified using \texttt{jupyter nbconvert --execute}
with a 120-second timeout per notebook. The full end-to-end retest across
all five notebook groups confirmed a 120/120 pass rate. Notebooks are
executed against the actual Zenodo corpus files to ensure that all
data access patterns are valid and all field names are correct.

%=====================================================================
\section{Planned LLM Fine-Tune Team}
\label{sec:llms}
%=====================================================================

Stage 2 of the platform (planned 2026) will add a four-model fine-tuned
LLM team covering hardware configurations from laptop to research cluster:

The platform's JSONL schema is explicitly designed for LLM RAG pipelines.
Stage~2 will provide a four-tier fine-tuned model team spanning
research-cluster to laptop hardware, covering parameter scales from
$\sim$72B (research-grade retrieval) to $\sim$7B (in-classroom,
CPU-compatible inference). All models will be fine-tuned on the
four EPS corpora using Unsloth QLoRA, with adapter weights published
to HuggingFace at \texttt{eps-research/} upon completion.
A vision-capable model tier ($\sim$72B) will support extraction of
kinematic values directly from ALMA plots and rotation curve figures
into structured JSONL.

%=====================================================================
\section{Data Access and Citation}
\label{sec:data}
%=====================================================================

All four corpora are released under CC BY 4.0 at Zenodo:

\begin{itemize}
\item Unified HI Corpus v7.0: \dataset[10.5281/zenodo.19563417]{https://doi.org/10.5281/zenodo.19563417}
\item Dwarf/Irregular Corpus v1.0: \dataset[10.5281/zenodo.20320362]{https://doi.org/10.5281/zenodo.20320362}
\item GC Corpus v1.3.1: \dataset[10.5281/zenodo.19907765]{https://doi.org/10.5281/zenodo.19907765}
\item High-z Corpus Z1: \dataset[10.5281/zenodo.20369285]{https://doi.org/10.5281/zenodo.20369285}
\end{itemize}

The platform software is archived at Zenodo
(DOI: \href{https://doi.org/10.5281/zenodo.20398430}{10.5281/zenodo.20398430})
and maintained at \url{https://github.com/eps-research/rag-corpus-series}.

Authors using any EPS Research corpus are requested to cite the relevant
Zenodo data descriptor and the Flynn \& Cannaliato (2025) omega correction
paper \citep{flynn2025} if the omega framework is used.

%=====================================================================
\section{Summary}
\label{sec:summary}
%=====================================================================

We have presented the EPS Research Astro-RAG Platform v1.0, providing:

\begin{enumerate}
\item Four open, machine-readable corpora totaling 772 objects spanning
      $z = 0$ to $z \sim 6$, with permanent Zenodo DOIs.
\item Full-fidelity preservation of the SPARC photometric decomposition
      database, including all baryonic components.
\item 120 verified executable Jupyter notebooks with a 120/120 end-to-end
      pass rate, covering rotation curve analysis, globular cluster dynamics,
      high-redshift kinematics, and educational outreach.
\item A QuickStart reproducibility pathway enabling independent verification
      of the cross-epoch omega sign reversal in under 10 minutes.
\item A High-School Exploration Track making real published astrophysical
      data accessible to students aged 12--18 without prior domain knowledge.
\item A documented roadmap for Stage 2 LLM fine-tune deployment on hardware
      ranging from free Google Colab to research-grade dual-GPU clusters.
\end{enumerate}

The cross-epoch omega sign reversal --- from $+7.06 \pm 3.26$ rad~Gyr$^{-1}$
at $z = 0$ to $-13.05$ rad~Gyr$^{-1}$ at $z \sim 5$ --- is the unifying
scientific result motivating the platform's design and its planned RAMSES
cosmological simulation program.

%=====================================================================
\section*{Declaration of Generative AI Use}
\label{sec:ai}

The author used AI language model assistants (Anthropic Claude,
Google Gemini) during the preparation of this manuscript for tasks
including code debugging, LaTeX formatting, and editorial review.
All scientific content, data analysis, corpus construction, and
conclusions are the sole responsibility of the author.

\begin{acknowledgments}
The author thanks J. Cannaliato for collaboration on the omega correction
framework. This research was conducted independently at EPS Research
without institutional affiliation or external funding. All data products
are released openly under CC BY 4.0.
\end{acknowledgments}

\facility{WSRT, VLA, ALMA, HST, Gaia}

\software{numpy \citep{harris2020}, matplotlib \citep{hunter2007},
          astropy \citep{astropy2022}, 3DBarolo \citep{diteodoro2015},
          jupyter \citep{kluyver2016}}

%=====================================================================
\begin{thebibliography}{}

\bibitem[Astropy Collaboration et al.(2022)]{astropy2022}
Astropy Collaboration et al. 2022, ApJ, 935, 167

\bibitem[Baumgardt \& Hilker(2018)]{baumgardt2018}
Baumgardt, H., \& Hilker, M. 2018, MNRAS, 478, 1520

\bibitem[Di Teodoro \& Fraternali(2015)]{diteodoro2015}
Di Teodoro, E.~M., \& Fraternali, F. 2015, MNRAS, 451, 3021

\bibitem[Flynn \& Cannaliato(2025)]{flynn2025}
Flynn, D.~C., \& Cannaliato, J. 2025, Frontiers in Astronomy and Space
Sciences, 12. \doi{10.3389/fspas.2025.1680387}

\bibitem[Flynn(2026a)]{flynn2026_hi}
Flynn, D.~C. 2026a, arXiv:2604.13489

\bibitem[Flynn(2026b)]{flynn2026_gc}
Flynn, D.~C. 2026b, arXiv:2605.03099

\bibitem[Flynn(2026c)]{flynn2026_dwarfs}
Flynn, D.~C. 2026c, arXiv:2605.22163

\bibitem[Flynn(2026d)]{flynn2026_z1}
Flynn, D.~C. 2026d, arXiv:2605.25339

\bibitem[Harris(1996)]{harris1996}
Harris, W.~E. 1996, AJ, 112, 1487

\bibitem[Harris et al.(2020)]{harris2020}
Harris, C.~R. et al. 2020, Nature, 585, 357

\bibitem[Hunter(2007)]{hunter2007}
Hunter, J.~D. 2007, Computing in Science and Engineering, 9, 90

\bibitem[Jones et al.(2021)]{jones2021}
Jones, G.~C. et al. 2021, MNRAS, 507, 3540

\bibitem[Kluyver et al.(2016)]{kluyver2016}
Kluyver, T. et al. 2016, in Positioning and Power in Academic Publishing,
87

\bibitem[Koribalski et al.(2019)]{koribalski2018}
Koribalski, B.~S. et al. 2019, MNRAS, 483, 4

\bibitem[Koribalski et al.(2020)]{koribalski2020}
Koribalski, B.~S. et al. 2020, Ap\&SS, 365, 118

\bibitem[Le F\`evre et al.(2020)]{lefevre2020}
Le F\`evre, O. et al. 2020, A\&A, 643, A1

\bibitem[Lelli et al.(2016)]{lelli2016}
Lelli, F., McGaugh, S.~S., \& Schombert, J.~M. 2016, AJ, 152, 157

\bibitem[Oh et al.(2015)]{oh2015}
Oh, S.-H. et al. 2015, AJ, 149, 180

\bibitem[Ott et al.(2012)]{ott2012}
Ott, J. et al. 2012, AJ, 144, 123

\bibitem[Vasiliev \& Baumgardt(2021)]{vasiliev2021}
Vasiliev, E., \& Baumgardt, H. 2021, MNRAS, 505, 5978

\bibitem[Walter et al.(2008)]{walter2008}
Walter, F. et al. 2008, AJ, 136, 2563

\end{thebibliography}

\end{document}
